\begin{frame}[plain]
\frametitle{Four cohorts, three prompt families}
\begin{tikzpicture}[remember picture,overlay]
  \sput{0.9}{1.50}{6.7}{\centering\fs{22}{26}\bfseries\color{SlideBlue}44,118}
  \sput{0.9}{2.35}{6.7}{\centering\fs{11}{14}Portuguese--English response pairs}
  \sput{8.4}{1.50}{6.7}{\centering\fs{22}{26}\bfseries\color{SlideBlue}14,706}
  \sput{8.4}{2.35}{6.7}{\centering\fs{11}{14}survey-specific records}
  \centerfigure{3.05}{\begin{tikzpicture}[text width=,anchor=center]
\path[use as bounding box] (0,0) rectangle (14.6,4.4);
\begin{axis}[
  slideaxis,width=7.5cm,height=4.3cm,
  at={(1.2cm,1.1cm)},anchor=south west,
  xbar,bar width=0.48cm,xmin=0,xmax=18000,
  ytick={0,1,2,3},yticklabels={A\quad2023,B\quad2024,C\quad2025,D\quad2025},
  y dir=reverse,ymin=-0.6,ymax=3.6,
  xtick={0,5000,10000,15000},xticklabels={0,5{,}000,10{,}000,15{,}000},
  xlabel={Response pairs},scaled x ticks=false,ymajorgrids=false,xmajorgrids=true,
  nodes near coords,
  every node near coord/.append style={font=\fs{9.5}{12},anchor=west,text=Ink,
    /pgf/number format/fixed,/pgf/number format/precision=0},
]
\addplot[fill=SlideBlue!72,draw=none]
  table[x=pairs,y expr=\coordindex,col sep=tab]{../shared/data/cohorts.tsv};
\end{axis}
\node[auditnode,text width=4.1cm,minimum height=1.1cm,inner sep=4pt] at (11.65,3.45)
  {Positive workplace qualities};
\node[auditnode,text width=4.1cm,minimum height=1.1cm,inner sep=4pt] at (11.65,2.0)
  {Areas for improvement};
\node[auditnode,text width=4.1cm,minimum height=1.1cm,inner sep=4pt] at (11.65,0.55)
  {Mental health and well-being};
\end{tikzpicture}
}
  \takeaway{Each retained record contributes one answer to each prompt family.}
  \source{Camera-ready article, Table 1 and cohort-selection description. Counts represent survey-specific records.}
\end{tikzpicture}
\note{
  - 01:35--02:35, 60 seconds\\
  - The corpus covers four Brazilian workplace survey cohorts from 2023 to 2025. There are 44,118 response pairs from 14,706 survey-specific records.\\
  - These records describe participation within each survey. A person answering across years cannot be linked across cohorts.\\
  - Every retained record has three answers, one per prompt family. One third of the pairs therefore come from mental-health or well-being prompts, whose wording varies by cohort.\\
  - The chart shows the cohort sizes. The analysis retained short and repeated answers because they are part of the operational setting.\\
  - Source: camera-ready article, Table 1 and Appendix A.
}
\end{frame}

\begin{frame}[plain]
\frametitle{Three audit views}
\begin{tikzpicture}[remember picture,overlay]
  \sput{1.15}{1.65}{14.6}{\begin{tikzpicture}[text width=,anchor=center]
\path[use as bounding box] (0,0) rectangle (14.6,5.3);
\node[auditnode,text width=2.2cm,minimum height=1.4cm] (pairs) at (1.6,2.65)
  {\textbf{Paired}\\responses};
\node[auditnode,text width=3.7cm,minimum height=1.05cm] (retrieval) at (6.4,4.35)
  {\textbf{Paired-text retrieval}\\{\fs{9}{12}Qwen3-Embedding-4B}};
\node[auditnode,text width=3.7cm,minimum height=1.05cm] (quality) at (6.4,2.65)
  {\textbf{Quality estimation}\\{\fs{9}{12}COMETKiwi DA-XL}};
\node[auditnode,text width=3.7cm,minimum height=1.05cm] (coverage) at (6.4,0.95)
  {\textbf{Coverage diagnostics}\\{\fs{9}{12}Lingua and EmoAtlas}};
\coordinate (branch) at (3.55,2.65);
\draw[MidGray,line width=0.8pt] (pairs.east) -- (branch);
\draw[auditline] (branch) |- (retrieval.west);
\draw[auditline] (branch) -- (quality.west);
\draw[auditline] (branch) |- (coverage.west);
\node[plainlabel,text width=4.5cm] at (11.7,4.35)
  {Can the paired text be recovered?};
\node[plainlabel,text width=4.5cm] at (11.7,2.65)
  {What quality does the model estimate?};
\node[plainlabel,text width=4.5cm] at (11.7,0.95)
  {Which language and emotion terms are recognized?};
\end{tikzpicture}
}
  \takeaway{Retrieval, quality estimates and coverage are interpreted separately.}
  \source{Camera-ready article, Figure 2 and Data and Methods. Model references appear on the final backup slide.}
\end{tikzpicture}
\note{
  - 02:35--03:40, 65 seconds\\
  - All pairs enter three parallel audit views. Read them as separate questions rather than as stages feeding a combined score.\\
  - Qwen3-Embedding-4B tests paired-text recovery. The 8B version provides a model-scale sensitivity check.\\
  - COMETKiwi DA-XL estimates quality from the Portuguese source and English translation, without a human reference translation.\\
  - Lingua supplies language flags, and EmoAtlas supplies language-specific affective-lexicon diagnostics. Coverage helps explain which responses yield a profile.\\
  - Source: camera-ready article, Figure 2 and Computational Provenance. Qwen3, COMETKiwi and EmoAtlas references are on the backup reference slide.
}
\end{frame}

\begin{frame}[plain]
\frametitle{What rank-one recovery measures}
\begin{tikzpicture}[remember picture,overlay]
  \centerfigure{1.45}{\begin{tikzpicture}[text width=,anchor=center]
\path[use as bounding box] (0,0) rectangle (14.6,4.15);
\node[auditnode,text width=2.5cm,minimum height=1.05cm] (query) at (1.9,2.85)
  {Portuguese\\query};
\node[plainlabel,text=SlideBlue,font=\fs{11}{14}\bfseries] at (9.5,3.85)
  {Ranked English candidates};
\node[auditnode,text width=5.7cm,minimum height=0.85cm,draw=SlideBlue]
  (first) at (9.5,2.85) {\textbf{Paired target or\\exact normalized duplicate}};
\node[auditnode,text width=5.7cm,minimum height=0.75cm,fill=white]
  at (9.5,1.50) {Different normalized wording};
\node[auditnode,text width=5.7cm,minimum height=0.75cm,fill=white]
  at (9.5,0.35) {Other candidate texts};
\draw[auditline] (query.east) -- (first.west);
\node[plainlabel,text=SlideBlue,font=\fs{10}{13}\bfseries] at (5.5,3.45) {Rank 1};
\node[plainlabel,text=MidGray,font=\fs{10}{13}] at (5.5,1.50) {Rank 2};
\node[plainlabel,text=MidGray,font=\fs{10}{13}] at (5.5,0.35) {Rank 3};
\end{tikzpicture}
}
  \sput{0.9}{5.95}{6.7}{\centering\fs{21}{25}\bfseries\color{SlideBlue}71.80\%}
  \sput{0.9}{6.65}{6.7}{\centering\fs{9.3}{11.5}English query\\to Portuguese candidates}
  \sput{8.4}{5.95}{6.7}{\centering\fs{21}{25}\bfseries\color{MidGray}73.37\%}
  \sput{8.4}{6.65}{6.7}{\centering\fs{9.3}{11.5}Portuguese query\\to English candidates}
  \takeaway{Every query is ranked against all 44,118 texts in the other language.}
  \source{Camera-ready article, retrieval definition and Table 2. Percentages are overall duplicate-aware R@1.}
\end{tikzpicture}
\note{
  - 03:40--04:45, 65 seconds\\
  - A query asks the encoder to find its counterpart among all 44,118 texts in the other language. Queries receive a retrieval instruction, while candidates are encoded without one.\\
  - Rank-one recovery counts a match when the first candidate is the paired target or any candidate with exactly the same normalized target text.\\
  - Semantically similar answers with different normalized wording remain competing candidates. This is a demanding corpus-level identification task.\\
  - The overall rates are 71.80 percent for English queries and 73.37 percent for Portuguese queries. The next slide shows what this average conceals.\\
  - Source: camera-ready article, Instruction-Conditioned Paired-Text Retrieval and Table 2. Diagram is schematic and contains no response text.
}
\end{frame}
